\section{Changes in version 3}\label{app:changes}

Version~2 of 27 September 2026 was withdrawn from the branch on the same day because of its framing, together with the audit notes. Version~3 keeps its theorems, proofs, citations and status lines, and makes the following changes.

\paragraph{Framing.}
\begin{itemize}
\item The front matter is a prose abstract and an introduction with the main results stated as Theorems~I--IV. The section \emph{About this reader} is replaced by the subsection \emph{How this record was made}, which gives the credit for the construction, the record and this paper, and the four statuses.
\item Remarks on how documents were produced, other than credit, are removed: the naming of the source on the monograph's title page, the correction note's author line, the record's self-description as a verification aid, the description of the record's labels as assigned by its verification, and the description of the conversation files. The statement that the source manuscript is unrefereed is removed.
\item The sentence that the paper takes no position on the construction is replaced by the statement that the construction and the counterexample claim are not verified here.
\item Section~\ref{sec:overlaps} is now \emph{Bridges to other areas and to the author's other workbenches}, and treats the connections to established theories first. The remark on the name of the construction is reduced to the statement that the $S^6$ construction and the Jacobian counterexample are different objects.
\item The paper is called a paper, and the companion texts are cited in their current versions.
\end{itemize}

\paragraph{New results.}
\begin{itemize}
\item Section~\ref{sec:monodromy}, the integral monodromy of the period system: Lemma~\ref{lem:stab}, Theorem~\ref{thm:mono}, Propositions~\ref{prop:heis} and~\ref{prop:twist}, Theorem~\ref{thm:periph}, Theorem~\ref{thm:orbitsS6} and Corollary~\ref{cor:orbitsS6}, and the proved negative statements of Section~\ref{ssec:mononeg}. They were first written in the audit note \note{52\_} and refereed there (Appendix~\ref{app:verification}).
\item Section~\ref{ssec:theories}: the paramodular group of level $6$ and the restrictions of $e(\Phi/12)$, spin structures, the Jordan type at the cusp, the modular part, and arithmeticity.
\item Questions~\ref{q:paramodular}--\ref{q:index12}.
\end{itemize}

\paragraph{Corrections.}
\begin{itemize}
\item The edge $S^6\to$ ES now states that the orbits of the covers at every level are determined (Theorem~\ref{thm:orbitsS6}); version~2 stated only that the package's statements are correct for odd levels and do not hold as stated at even levels, which remains true.
\end{itemize}
